\documentclass[9pt]{extarticle}
\usepackage[a4paper,margin=0.25in]{geometry}
\usepackage{lmodern}
\usepackage{multicol}
\usepackage{amsmath,amssymb}
\usepackage{graphicx}
\usepackage{eso-pic}
\usepackage{enumitem}
\usepackage{titlesec}
\usepackage[dvipsnames]{xcolor}
\usepackage{tcolorbox}
\tcbuselibrary{skins,breakable}
\usepackage{fancyhdr}
\pagestyle{fancy}
\fancyhf{}
\renewcommand{\headrulewidth}{0pt}
\setlength{\parindent}{0pt}
\setlength{\columnsep}{0.15in}
\setlength{\columnseprule}{0.4pt}
\setlength{\abovedisplayskip}{2pt}
\setlength{\belowdisplayskip}{2pt}
\setlength{\abovedisplayshortskip}{1pt}
\setlength{\belowdisplayshortskip}{1pt}
\allowdisplaybreaks
\emergencystretch=2em

% Compact list environment
\setlist[itemize]{itemsep=0pt,parsep=0pt,topsep=1pt,leftmargin=*,partopsep=0pt}
\setlist[enumerate]{itemsep=0pt,parsep=0pt,topsep=1pt,leftmargin=*,partopsep=0pt}

\titleformat{\section}{\normalfont\bfseries\small\color{RoyalBlue}}{\thesection}{0.3em}{{\color{black!60}\titlerule[0.9pt]}}
\titlespacing{\section}{0pt}{2pt}{1pt}
\titleformat{\subsection}{\normalfont\bfseries\footnotesize\color{OliveGreen}}{}{0em}{\titlerule[0.4pt]}
\titlespacing{\subsection}{0pt}{1.5pt}{1pt}
\setlength{\parskip}{0pt}
\linespread{1.2}

\newcommand{\hl}[1]{\textbf{\textcolor{BrickRed}{#1}}}
\newcommand{\wk}[1]{\textbf{\textcolor{OliveGreen}{#1}}}
\newcommand{\ex}[1]{\textbf{\textcolor{RoyalBlue}{#1}}}
\newcommand{\sepline}{\vspace{1pt}{\color{RoyalBlue}\hrule}\vspace{3pt}}

\newtcolorbox{kbox}{colback=Goldenrod!10,colframe=RoyalBlue!70,boxrule=0.6pt,arc=1.5pt,left=1.5pt,right=1.5pt,top=0.5pt,bottom=0.5pt,boxsep=0pt,before skip=1pt,after skip=1pt}

\begin{document}
\footnotesize

\AddToShipoutPicture*{%
\AtTextUpperLeft{%
\raisebox{6pt}{\makebox[\textwidth][r]{\normalfont\scriptsize\textcolor{black!60}{Updated: 23 Sep 2026 19:34}}}%
}%
}
\begin{center}
\colorbox{ForestGreen!12}{\parbox{0.97\linewidth}{\centering
{\bfseries\large\color{ForestGreen} Numerical Analysis (SCI19 3111) --- Final Cheat Sheet}\\[1.5pt]
{\scriptsize\color{ForestGreen!70!black} Regression $\cdot$ Integration $\cdot$ ODEs $\cdot$ Linear Algebra (+ midterm recap)}}}
\end{center}
\sepline

\begin{multicols}{3}

\section{Regression / Data Fitting (Th.1--6)}
\begin{sloppypar}\raggedright
\textbf{Regression $\ne$ interpolation:} interpolation trusts every measurement and passes through ALL points; regression assumes noise and finds the best OVERALL trend, passing through (almost) none.
Deviations $d_i=y_i-\phi(x_i)$, data $(x_i,y_i)$, $i=0..n$ ($n{+}1$ points).
\textbf{Error choices:} naive $S=\sum d_i$ (rejected: $+,-$ cancel); abs $S=\sum|d_i|$ (rejected: kink at 0, no derivative); \hl{squared} $S=\sum d_i^2$ (chosen: smooth everywhere, so ``set gradient $=0$'' gives solvable linear equations; big misses punished quadratically).
Example (Lec14): pts $(1,2),(2,1)$, line through $(1,1),(2,2)$: $d_0=1,d_1=-1$: naive $S=0$ (cancels!), abs $S=2$, squared $S=2$.

\subsection{Linear regression}
$\hat y=a_0+a_1x$, $S=\sum(y_i-a_0-a_1x_i)^2$. General fit $\phi(x,a_0{..}a_\ell)$: $\tfrac{\partial S}{\partial a_k}=-2\sum_{i=0}^n(y_i-\phi(\cdots))\tfrac{\partial\phi}{\partial a_k}=0$, $k=0..\ell$ ($\ell{+}1$ equations). For the line this gives the two normal equations below.
$\partial_{a_0}$: $\sum y_i-(n{+}1)a_0-a_1\sum x_i=0$; $\partial_{a_1}$: $\sum x_iy_i-a_0\sum x_i-a_1\sum x_i^2=0$.\\
Divide by $n{+}1$ $\Rightarrow$ $a_0+a_1\bar x=\bar y$, $a_0\bar x+a_1\overline{x^2}=\overline{xy}$; Cramer gives
\begin{kbox}
\centering
$\displaystyle a_1=\frac{\overline{xy}-\bar x\bar y}{\overline{x^2}-\bar x\cdot\bar x},\qquad a_0=\bar y-a_1\bar x$
\end{kbox}
with $\bar x=\tfrac1{n+1}\sum x_i$, $\bar y=\tfrac1{n+1}\sum y_i$, $\overline{x^2}=\tfrac1{n+1}\sum x_i^2$, $\overline{xy}=\tfrac1{n+1}\sum x_iy_i$.\\
\hl{Trap:} $\overline{x^2}\ne\bar x^2$; denominator $=0$ iff all $x_i$ equal (no fit possible).

\subsection{Hessian + minimum test}
\begin{kbox}
\centering
$H=\left(\begin{smallmatrix}S_{a_0a_0}&S_{a_0a_1}\\S_{a_1a_0}&S_{a_1a_1}\end{smallmatrix}\right)$\\[-2pt]
$\displaystyle S_{a_0a_0}=2(n{+}1),\ S_{a_0a_1}=2\sum x_i,\ S_{a_1a_1}=2\sum x_i^2$\\[1pt]
$\displaystyle \det H=4\big[(n{+}1)\sum x_i^2-(\sum x_i)^2\big]$; min $\iff S_{a_0a_0}>0$ AND $\det H>0$
\end{kbox}
Intuition: $H$ = error-bowl curvature; \textbf{Cauchy--Schwarz} $\Rightarrow\det H\ge0$, $>0$ if $\ge2$ distinct $x_i$ $\Rightarrow$ unique global min.

\subsection{Multivariable + pseudoinverse}
$\hat y=a_0+a_1x_1+\cdots+a_mx_m$, $X$ is $(n{+}1)\times(m{+}1)$ with first col of 1's. $S=\|\bar y-X\bar a\|^2$.
\begin{kbox}
\centering
$\displaystyle X^TX\bar a=X^T\bar y,\qquad \bar a=\underbrace{(X^TX)^{-1}X^T}_{X^+}\bar y$
\end{kbox}
$X^+=$ Moore--Penrose pseudoinverse ($=X^{-1}$ when $X$ square). Idea: projects $\bar y$ onto the column space of $X$. \hl{Trap:} $X^+X=I$ holds; $XX^+=I$ generally does NOT.

\subsection{ML: classification}
Features = inputs $(x_1,x_2)$; plane $a_0+a_1x_1+a_2x_2=0\cap x_1x_2$-plane = \hl{decision boundary} (class 1 vs 2). Logistic $\sigma(z)=\tfrac1{1+e^{-z}}: \mathbb R\to[0,1]$; classifier $\sigma(a_0+a_1x_1+a_2x_2)$. \ex{Worked mini-example:} plane $-1+x_1+x_2=0$: $(0,0)\to-1<0$ (class 1), $(1,1)\to+1>0$ (class 2); $\sigma$ gives $\approx0.27$, $\approx0.73$.
\end{sloppypar}

\section{Linearization}
\wk{Goal:} transform $(\tilde x,\tilde y)$ so model is linear; fit; inverse-transform back. Logs turn $\times$ into $+$ and powers into $\times$. Steps: (1) change vars; (2) fit $\tilde y=a_0+a_1\tilde x$; (3) recover $A,B$; (4) check fit.
\begin{itemize}
\item $y=Ae^{Bx}$: $\tilde x=x$, $\tilde y=\ln y$; $a_0=\ln A$, $a_1=B$ $\Leftrightarrow$ $\textcolor{BrickRed}{A=e^{a_0}}$.
\item $y=Ax^n$: $\tilde x=\ln x$, $\tilde y=\ln y$; $a_0=\ln A$, $a_1=n$ $\Leftrightarrow$ $\textcolor{BrickRed}{A=e^{a_0}}$.
\item $y=A+\tfrac{B}{x}$: $\tilde x=\tfrac1x$, $\tilde y=y$; $a_0=A$, $a_1=B$.
\item $y=\tfrac1{A+Bx}$: $\tilde x=x$, $\tilde y=\tfrac1y$; $a_0=A$, $a_1=B$.
\item $y=\tfrac{Ax}{B+x}$: $\tilde x=\tfrac1x$, $\tilde y=\tfrac1y$; $a_0=\tfrac1A$, $a_1=\tfrac{B}{A}$ $\Leftrightarrow$ $\textcolor{BrickRed}{A=\tfrac1{a_0},\,B=\tfrac{a_1}{a_0}}$.
\end{itemize}
\hl{Trap:} not everything linearizes (e.g. $a_0e^{-x}\sin(a_1+2\pi x)$). \wk{Model from graph:} straight $\to$ linear; explosive growth $\to Ae^{Bx}$; power curve $\to Ax^n$; saturating hyperbola $\to Ax/(B+x)$; reciprocal decay $\to A+B/x$.

\section{Grids + Riemann Sums}
\begin{sloppypar}\raggedright
Partition $a=x_0<\cdots<x_n=b$, $h_i=x_i-x_{i-1}$, sample $x_i^*\in[x_{i-1},x_i]$: $\int_a^bf\approx\sum f(x_i^*)h_i$.
\textbf{Uniform:} $h_i\equiv h=\tfrac{b-a}{n}$, $n{+}1$ points. \textbf{Non-uniform:} varying $h_i$ (adaptive: fine where $f$ changes fast, coarse elsewhere).
\textbf{Riemann integral} $=\lim_{\max h_j\to0}\sum h_if(x_i^*)$; $f$ bounded, at most countably many jumps.
\end{sloppypar}

\section{Rectangle / Trapezoid / Simpson}
\begin{sloppypar}\raggedright
\textbf{Sample point:} $x_i^*=\alpha x_{i-1}+(1-\alpha)x_i$: $\alpha=1$ left, $0$ right, $\tfrac12$ mid. (e.g. Quiz7 $\alpha=0.2$.)
General non-uniform: left $\sum h_if(x_{i-1})$, right $\sum h_if(x_i)$, mid $\sum h_if(\tfrac{x_{i-1}+x_i}2)$.
\begin{kbox}
\centering
$\tilde I_{\text{Tr}}=\tfrac12\sum_{i=1}^n\big(f(x_{i-1})+f(x_i)\big)h_i$
\end{kbox}
Uniform trapezoid: $\tfrac{h(f_0+f_n)}2+h\sum_{i=1}^{n-1}f_i$.
\wk{Trap idea:} one strip = rectangle $hQ_1$ + triangle $\tfrac h2(Q_2-Q_1)$ $=h(Q_1+Q_2)/2$; uniform version averages neighbors (interior $f_i$ counted twice $\to$ single sum).
\wk{Iterative (halve $h$, reuse):} $S_1=f_0+f_1$; double $N$ by adding $2\times$(new-node $f$'s): $S_2=S_1+2f_1^{\text{new}}$, $S_4=S_2+2(f_1+f_3)^{\text{new}}$ — evaluate only new nodes (relabelled each doubling).
\textbf{Simpson} ($n=2m$ even, uniform): fits parabola on each pair (one pair needs 2 subintervals $\Rightarrow$ $n$ even). Why $O(h^4)$? Parabola matches value+slope+curvature (trap: only value+slope).
\begin{kbox}
\centering
$I_{\text{parab}}=\tfrac{h}{3}\big\{f_0+f_n+4\sum_{\text{odd }i}f_i+2\sum_{\text{even }i<n}f_i\big\}$
\end{kbox}

\subsection{Accuracy estimates}
\begin{itemize}
\item left/right rectangle $O(h)$; midpoint $O(h^2)$; trapezoid $O(h^2)$; Simpson $O(h^4)$.
\item Read $O(h^p)$ as: halving $h$ divides the error by $2^p$ (rect $\div2$, mid/trap $\div4$, Simpson $\div16$). E.g. $\int_0^2x^3dx=4$: $T_1=8,T_2=5,T_4=4.25$; err $4,1,0.25$ ($\div4$ each halving).
\item Left-rect global: $|E|\le\tfrac{M_1(b-a)}2h$, $M_1=\max|f'|$. Midpoint: $|E|\le\tfrac{M_2(b-a)}{24}h^2$ (half the trapezoid constant $1/12$). (Taylor + symmetry kills $f'$ term.)
\item No elementary antiderivative ($\int e^{x^2}dx$, $\int\tfrac{dx}{\ln x}$): cross-check Simpson vs trapezoid.
\end{itemize}
\end{sloppypar}

\section{ODEs: IVP + Numeric Methods}
\begin{sloppypar}\raggedright
$y'=f(x,y)$ (slope field = direction at each point). General solution $y(x,C)$ = whole family of curves; IVP (Cauchy) $+y(x_0)=y_0$ picks the ONE curve through $(x_0,y_0)$ (fix $C$ by plugging IC; existence/uniqueness assumed).
\hl{Check-solution trap:} differentiate candidate AND plug into ODE — both must hold for ALL $C$.
\wk{Bridge (Lec19):} Euler comes from the finite difference $y'\approx\tfrac{y_{i+1}-y_i}{h}$: replace $y'$ in $y'=f(x,y)$ and solve for $y_{i+1}$.

\subsection{Finite differences (VII-1)}
\textbf{fwd:} $\tfrac{f(x+h)-f(x)}h$ $O(h)$; \textbf{bwd:} $\tfrac{f(x)-f(x-h)}h$ $O(h)$; \textbf{central:} $\tfrac{f(x+h)-f(x-h)}{2h}$ $O(h^2)$ = avg of fwd+bwd. \textbf{2nd:} central $f''\approx\tfrac{f(x+h)-2f(x)+f(x-h)}{h^2}$ $O(h^2)$. \textbf{Best $h$:} $(3\epsilon/M)^{1/3}$ ($\epsilon$ = rounding err, $M=\max|f'''|$).
\begin{kbox}
\centering
$\displaystyle y_{i+1}=y_i+hf(x_i,y_i),\quad x_{i+1}=x_i+h\quad\text{(Euler)}$
\end{kbox}
\begin{kbox}
\centering
$\displaystyle y_{i+1}=y_i+hf(x_{i+\frac12},y_{i+\frac12}),\quad x_{i+\frac12}=x_i+\tfrac h2,\ y_{i+\frac12}=y_i+\tfrac h2f(x_i,y_i)$
\end{kbox}
\textbf{RK4:} weights $(1,2,2,1)/6$: $k_1=f(x_n,y_n)$; $k_2=f(x_n{+}\tfrac h2,y_n{+}\tfrac h2k_1)$; $k_3=f(x_n{+}\tfrac h2,y_n{+}\tfrac h2k_2)$; $k_4=f(x_n{+}h,y_n{+}hk_3)$.
\begin{kbox}
\centering
$\displaystyle \textcolor{BrickRed}{y_{n+1}=y_n+\tfrac h6(k_1{+}2k_2{+}2k_3{+}k_4)}$
\end{kbox}
\ex{Write-only schemes (Quiz9 Q4 + HW10 Q3):} for $y'=y+x$: $k_1=y_i{+}x_i$, $k_2=y_i{+}\tfrac h2k_1{+}x_i{+}\tfrac h2$, $k_3=y_i{+}\tfrac h2k_2{+}x_i{+}\tfrac h2$, $k_4=y_i{+}hk_3{+}x_i{+}h$. For $y'=xy+1/x$ ($h{=}1/12$): $k_1=x_iy_i+1/x_i$; $k_2=(x_i{+}\tfrac1{24})(y_i{+}\tfrac{k_1}{24})+1/(x_i{+}\tfrac1{24})$; $k_3$ same with $k_2$; $k_4=(x_i{+}\tfrac1{12})(y_i{+}\tfrac{k_3}{12})+1/(x_i{+}\tfrac1{12})$.
\textbf{Global errors:} Euler $O(h)$ (local $O(h^2)\times N\sim1/h$); midpoint $O(h^2)$ (local $O(h^3)$, symmetric); RK4 $O(h^4)$ (4 evals/step).
\end{sloppypar}

\section{Linear Algebra}
\begin{sloppypar}\raggedright
$\det=$ oriented area (2D) / volume (3D), up to sign; $\det\ne0\iff$ invertible. $\det(AB)=\det A\det B$. $Ax=b$ unique $\iff\det A\ne0\iff n$ pivots.\\
\textbf{Scalar product:} $\bar a{\cdot}\bar b=a^xb^x+a^yb^y=|\bar a||\bar b|\cos\alpha$ (+$a^zb^z$ in 3D); $|\bar a|=\sqrt{(a^x)^2+(a^y)^2}$.\\
\ex{Basic ops:} $\big(\begin{smallmatrix}1&2\\3&4\end{smallmatrix}\big)\big(\begin{smallmatrix}5\\6\end{smallmatrix}\big)=\big(\begin{smallmatrix}17\\39\end{smallmatrix}\big)$ (row $\cdot$ column).\\
\textbf{Projection:} $\mathrm{proj}_{\bar b}\bar a=\tfrac{\bar a\cdot\bar b}{|\bar b|}=|\bar a|\cos\alpha$; basis: $\mathrm{proj}_{\bar\imath}\bar a=a^x$, $\mathrm{proj}_{\bar\jmath}\bar a=a^y$ (\ex{e.g.} $(0,1),(2,0)$: $1,2$).\\
\textbf{Area:} parallelogram $S=a^yb^x-a^xb^y=-\det$; area $=|S|=|\det|$. \ex{Worked:} $u=(3,1),v=(1,2)$: $S=1\cdot1-3\cdot2=-5$, area $5$. \textbf{Volume (3D):} $|\det|$ of the $3\times3$; \ex{Worked:} $\mathrm{diag}(1,2,3)$ volume $=6$. \ex{Non-diagonal:} sheet's $3\times3$ $A$ has $\det=-3\Rightarrow$ volume $3$.\\
\textbf{Minor/cofactor:} $\mathrm{Minor}(a_{ij})$ = det deleting row $i$, col $j$; $\mathrm{cofac}(a_{ij})=(-1)^{i+j}\mathrm{Minor}(a_{ij})$. \ex{Worked:} $\mathrm{cofac}(a_{23})=-(a_{11}a_{32}-a_{31}a_{12})=a_{31}a_{12}-a_{11}a_{32}$.\\
\textbf{Laplace:} $\det A=\sum_j a_{kj}\mathrm{cofac}(a_{kj})$ (any row $k$). \textbf{Row swap} flips sign. \ex{Worked:} $A=\big(\begin{smallmatrix}1&2&3\\4&5&6\\7&8&10\end{smallmatrix}\big)$: cofactors $2,2,-3$ $\Rightarrow$ $\det=1(2)+2(2)+3(-3)=-3$.\\
\textbf{Triangular:} zeros above (lower) or below (upper) diagonal; $\det=$ product of diagonal. \ex{Worked (5$\times$5):} $\mathrm{diag}(1,2,3,4,5)$ $\Rightarrow$ $\det=120=5!$.\\
\textbf{LU:} $A=LU$ ($L$ unit-lower, $U$ upper; set $\ell_{ii}=1$). \ex{Worked (2$\times$2):} $\big(\begin{smallmatrix}1&2\\3&4\end{smallmatrix}\big)=\big(\begin{smallmatrix}1&0\\3&1\end{smallmatrix}\big)\big(\begin{smallmatrix}1&2\\0&-2\end{smallmatrix}\big)$ ($u_{22}=4-3\cdot2=-2$). \textbf{Det:} $\det A=\det L\cdot\det U$; here $1\cdot(-2)=-2$ \wk{Check} $=1\cdot4-2\cdot3$.\\
\ex{Worked (3$\times$3):} $\big(\begin{smallmatrix}1&2&3\\4&5&6\\7&8&10\end{smallmatrix}\big)=\big(\begin{smallmatrix}1&0&0\\4&1&0\\7&2&1\end{smallmatrix}\big)\big(\begin{smallmatrix}1&2&3\\0&-3&-6\\0&0&1\end{smallmatrix}\big)$ ($\ell_1{=}4,u_4{=}-3,u_5{=}-6,\ell_2{=}7,\ell_3{=}2,u_6{=}1$); $\det=1\cdot(-3)\cdot1=-3$.\\
\textbf{Solve $Ax=y$:} $Ax=L(Ux)=y$; set $z=Ux$: solve $Lz=y$ (forward), then $Ux=z$ (backward). \ex{Worked:} $A$ above, $y=[1,-1]^T$: $z_1=1$, $3+z_2=-1\Rightarrow z_2=-4$; $-2x_2=-4\Rightarrow x_2=2$, $x_1+4=1\Rightarrow x_1=-3$. $\boxed{x=[-3,2]^T}$ \wk{Check:} $Ax=[1,-1]^T$.
\end{sloppypar}

\section{Midterm Recap (compact)}
\begin{sloppypar}\raggedright
\begin{itemize}
\item \textbf{Taylor:} $f(x)=\sum_{k=0}^n\tfrac{f^{(k)}(x_0)}{k!}(x-x_0)^k+R_n$, $|R_n|\le\tfrac{Md^{n+1}}{(n+1)!}$; odd fn $\to$ odd powers only (deg 8 $\to x^7$), bound $\tfrac{12^{2m+2}}{(2m+2)!}$. \textbf{Maclaurin:} $e^x,\sin,\cos$ $\forall\mathbb R$; $\tfrac1{1-x},\ln(1{+}x),\arctan$: $|x|<1$.
\item \textbf{Horner:} tableau at $c$ $\to$ $p(c)$ + quotient $q$; again on $q$ $\to$ $p'(c)$. \textbf{Interp:} $\ell_i=\prod_{j\ne i}\tfrac{x-x_j}{x_i-x_j}$, $p=\sum y_i\ell_i$ (err $0$ at nodes); Vandermonde $\det=\prod_{i<j}(x_j-x_i)\ne0$; DD $f[x_i,x_{i+1}]=\tfrac{\Delta f}{\Delta x}$ (Newton not required).
\end{itemize}
\end{sloppypar}

\section{Worked Mini-Examples}
\begin{sloppypar}\raggedright

\subsection{Coffee regression (Quiz 6)}
Given $x:0,1,2,3$; $y:3,5,8,7$ ($n{+}1{=}4$ pts).
\wk{Step 1} — sums (add term by term): $\sum x_i=0{+}1{+}2{+}3=6$; $\sum y_i=3{+}5{+}8{+}7=23$; $\sum x_i^2=0{+}1{+}4{+}9=14$; $\sum x_iy_i=0{+}5{+}16{+}21=42$.
\wk{Step 2} — means (divide by 4): $\bar x=6/4=1.5$; $\bar y=23/4=5.75$; $\overline{x^2}=14/4=3.5$; $\overline{xy}=42/4=10.5$.
\wk{Step 3} — coeffs (plug into boxed formula): $a_1=\tfrac{10.5-1.5\cdot5.75}{3.5-1.5^2}$\\
$=\tfrac{10.5-8.625}{3.5-2.25}=\tfrac{1.875}{1.25}=1.5$; $a_0=5.75-1.5\cdot1.5$\\
$=5.75-2.25=3.5$. $\boxed{\hat y=3.5+1.5x}$.
\wk{Step 4} — predict $x{=}4$: $\hat y=3.5+1.5\cdot4=3.5+6=9.5$ (mark $(4,9.5)$ on the scatter plot with the line).
\wk{Plot recipe (Quiz6 Q2):} scatter the 4 data points, draw the line through $(0,3.5)$ and $(4,9.5)$ — regression line need NOT pass through data points.
\wk{Step 5} — \wk{Check fit:} $\hat y:3.5,5.0,6.5,8.0$ vs $y:3,5,8,7$; residuals $-0.5,0,1.5,-1.0$; squared error $S=0.25+0+2.25+1=3.5$ (this $S$ is what least squares minimized).
\ex{Hessian (Quiz6 Q4, coffee data):} from Step 1 sums ($\sum x_i=6$, $\sum x_i^2=14$): $S_{a_0a_0}=2\cdot4=8$, $S_{a_0a_1}=2\cdot6=12$, $S_{a_1a_1}=2\cdot14=28$;\\
$\det H=8\cdot28-12^2=224-144=80>0$ (+ $8>0$) $\Rightarrow$ min.

\subsection{Battery fit (Lec 14)}
Given $t:0,10,20,30,40,50,60$; charge $5,10,15,25,35,40,45$ ($n{+}1{=}7$).
Sums: $\sum t_i=0{+}10{+}20{+}30{+}40{+}50{+}60=210$; $\sum y_i=5{+}10{+}15{+}25{+}35{+}40{+}45=175$; $\sum t_i^2=0{+}100{+}400{+}900{+}1600{+}2500{+}3600=9100$; $\sum t_iy_i=0{+}100{+}300{+}750{+}1400{+}2000{+}2700=7250$.
Means (divide by 7): $\bar x=210/7=30$; $\bar y=175/7=25$; $\overline{x^2}=9100/7=1300$; $\overline{xy}=7250/7\approx1035.71$.
Coeffs: $a_1=\tfrac{1035.71-30\cdot25}{1300-30^2}=\tfrac{1035.71-750}{1300-900}$\\
$=\tfrac{285.71}{400}\approx0.714$; $a_0=25-0.714285\cdot30$\\
$=25-21.43\approx3.57$. $\boxed{\phi(x)=0.714x+3.57}$.
\wk{Check fit:} $\hat y:3.57,10.71,17.85,24.99,32.13,39.27,46.41$ vs $y:5,10,15,25,35,40,45$.
\ex{Predict $t{=}70$ (same battery data):} $0.714\cdot70+3.57=53.55\%$ (extrapolation — flag as less reliable).

\subsection{Linearization drills (HW09 Q4 + Lec15)}
Given $x:1,2,3,5$, $y:8,4,2,1$; model $y=Ae^{Bx}$ $\Rightarrow$ transform $\tilde x=x$, $\tilde y=\ln y$.
\wk{Step 0} — transform each $y$: $\ln8\approx2.07944$, $\ln4\approx1.38629$, $\ln2\approx0.69315$, $\ln1=0$. So $\tilde x:1,2,3,5$; $\tilde y:2.07944,1.38629,0.69315,0$.
\wk{Step 1} — means ($n{+}1{=}4$): $\bar{\tilde x}=(1{+}2{+}3{+}5)/4=11/4=2.75$; $\bar{\tilde y}=(2.07944{+}1.38629{+}0.69315{+}0)/4=4.15888/4\approx1.03972$; $\overline{\tilde x^2}=(1{+}4{+}9{+}25)/4=39/4=9.75$; $\overline{\tilde x\tilde y}=(1\cdot2.07944{+}2\cdot1.38629{+}3\cdot0.69315{+}5\cdot0)/4=6.93147/4\approx1.73287$.
\wk{Step 2} — fit: $a_1=\tfrac{1.73287-2.75\cdot1.03972}{9.75-2.75^2}=\tfrac{1.73287-2.85923}{9.75-7.5625}$\\
$=\tfrac{-1.12636}{2.1875}\approx-0.5149$; $a_0=1.03972-(-0.5149)\cdot2.75$\\
$=1.03972+1.41598\approx2.4557$.
\wk{Step 3} — back-transform ($A=e^{a_0}$, $B=a_1$): $A=e^{2.4557}\approx11.65$, $B\approx-0.515<0$. \wk{Check:} $\hat y=11.65e^{-0.515x}$ gives $7.0,4.2,2.5,0.9$ vs $8,4,2,1$ — reasonable.\\
\ex{Second model (Lec15):} $x:1,2,3,4$, $y:0.45,0.81,1.05,1.13$ for $y=\tfrac{Ax}{B+x}$: $\tilde x=1/x$, $\tilde y=1/y$ give $a_1\approx1.834$, $a_0\approx0.368$,\\
so $A=1/a_0\approx2.71$, $B=a_1/a_0\approx4.98$.\\
\ex{Third model — $y=A+B/x$:} $x:1,2,4$, $y:5,3,2$: $\tilde x=1/x:1,0.5,0.25$; means $\bar{\tilde x}=0.58333$, $\bar y=3.33333$, $\overline{\tilde x^2}=0.4375$, $\overline{\tilde x y}=2.33333$; $a_1=B=\tfrac{2.33333-0.58333\cdot3.33333}{0.4375-0.58333^2}$\\
$=\tfrac{0.38889}{0.09722}=4.0$; $a_0=A=3.33333-4\cdot0.58333=1.0$. $\boxed{y=1+4/x}$ (exact through all 3 pts here).

\subsection{Trapezoid $\sin x$ (Quiz 7)}
Given $[0,3\pi/2]$, uniform $h=\pi/4$ $\Rightarrow$ $n=(3\pi/2)/(\pi/4)=6$ subintervals, 7 nodes $x_i=i\pi/4$.
\wk{Step 1} — evaluate $f(x)=\sin x$: $\sin0=0$; $\sin\frac\pi4=\tfrac{\sqrt2}2\approx0.70711$; $\sin\frac\pi2=1$; $\sin\frac{3\pi}4=\tfrac{\sqrt2}2$; $\sin\pi=0$; $\sin\frac{5\pi}4=-\tfrac{\sqrt2}2$; $\sin\frac{3\pi}2=-1$.
\wk{Step 2} — uniform trapezoid (endpoints half-weighted): $\tilde I=\tfrac h2(f_0{+}f_6)+h\sum_{i=1}^{5}f_i$ with $\sum_{1}^{5}f_i=\tfrac{\sqrt2}2+1+\tfrac{\sqrt2}2+0-\tfrac{\sqrt2}2=1+\tfrac{\sqrt2}2\approx1.70711$.
\wk{Step 3} — multiply out: $\tilde I=\tfrac h2(f_0{+}f_6)+h\sum_{1}^{5}f_i$\\
$=\tfrac\pi4\big(\tfrac12(-1)+1.70711\big)=\tfrac\pi4(1.20711)\approx0.78540\cdot1.20711\approx0.94806$.
\wk{Step 4} — exact $I=\int_0^{3\pi/2}\sin x\,dx=[-\cos x]_0^{3\pi/2}=-\cos\tfrac{3\pi}2+\cos0=0+1=1$; err $=|1-0.94806|\approx0.05194$.
\wk{Draw (Pract.6):} mark $n{+}1$ nodes; plot $(x_i,f(x_i))$; join with segments (trapezoids); Simpson: one arc per pair through its 3 points.

\subsection{Simpson $\int_2^4\tfrac{dx}{\ln x}$ (Quiz 8)}
Given $n=8$ (even $\checkmark$), 9 points, $h=(4-2)/8=0.25$. Nodes $x_i=2+0.25i$: $2,2.25,2.5,2.75,3,3.25,3.5,3.75,4$.
\wk{Step 1} — evaluate $f(x)=1/\ln x$ at each node: $f_0=1/\ln2\approx1.44270$; $f_1=1/\ln2.25\approx1.23315$; $f_2=1/\ln2.5\approx1.09136$; $f_3=1/\ln2.75\approx0.98853$; $f_4=1/\ln3\approx0.91024$; $f_5=1/\ln3.25\approx0.84842$; $f_6=1/\ln3.5\approx0.79824$; $f_7=1/\ln3.75\approx0.75657$; $f_8=1/\ln4\approx0.72135$.
\wk{Step 2} — weight (ends $\times1$, odd $i$ $\times4$, even $i$ $\times2$): $4\sum_{\text{odd}}=4(1.23315{+}0.98853{+}0.84842{+}0.75657)=4(3.82667)=15.30668$; $2\sum_{\text{even}}=2(1.09136{+}0.91024{+}0.79824)=2(2.79984)=5.59968$.
\wk{Step 3} — assemble: $I=\tfrac{0.25}3(1.44270+15.30668+5.59968+0.72135)$\\
$=\tfrac{0.25}3(23.07041)\approx0.083333\cdot23.07041\approx1.92253$.

\subsection{Euler + midpoint $y'=x-4y$, $y(0)=-3/4$ (Quiz 9)}
$h=0.2$, $f(x,y)=x-4y$. \ex{Euler:} $f_0=0-4(-0.75)=3$: $y_1=-0.75+0.2(3)=-0.15$; $f_1=0.2-4(-0.15)=0.8$: $y_2=-0.15+0.2(0.8)=0.01$; $f_2=0.4-0.04=0.36$: $y_3=0.01+0.2(0.36)=0.082$.
\ex{Midpoint:} half-step $(0.1,-0.75+0.1(3))=(0.1,-0.45)$, $f=0.1+1.8=1.9$ $\to y_1=-0.75+0.2(1.9)=-0.37$; half-step $(0.3,-0.37+0.1(1.68))=(0.3,-0.202)$, $f=0.3+0.808=1.108$ $\to y_2=-0.37+0.2(1.108)=-0.1484$.
Exact $y(0.4)\approx-0.101304$ (given; $y=x/4-1/16-\tfrac{11}{16}e^{-4x}$): Euler err $=|{-0.101304}-0.01|\approx0.11130$; midpoint err $=|{-0.101304}-(-0.1484)|\approx0.04710$.
\ex{Quiz9 Q1 — IVP from general sol:} $y=-\tfrac14+\tfrac x2+Ce^{-2x}$, $(x_0,y_0)=(\tfrac12,\tfrac14)$: $\tfrac14=-\tfrac14+\tfrac14+Ce^{-1}$ $\Rightarrow$ $C=e/4\approx0.6796$, $y^*=-\tfrac14+\tfrac x2+\tfrac e4e^{-2x}$.
\ex{HW10 Q2 pattern:} $y'=\tfrac yx+2x^2$, $(1,4)$, general $y=(x^2+C)x$: $4=(1+C)\cdot1\Rightarrow C=3$, $y^*=(x^2+3)x$. Euler pattern: $x_1=1+h$, $y_1=4+h(4/1+2)=4+6h$ (since $f(1,4)=6$). Midpoint pattern: $x_{1/2}=1+h/2$, $y_{1/2}=4+(h/2)\cdot6=4+3h$.
\hl{Quiz9 Q2 — check trap (fail case):} $y=Ce^x-2\cos2x-2\sin2x$ is \emph{not} a sol of $y'=y+5\sin2x$: $y'=Ce^x+4\sin2x-4\cos2x$, but $y+5\sin2x=Ce^x-2\cos2x+3\sin2x$; difference $=\sin2x-2\cos2x\ne0$.\\
\ex{Pass case (HW10 Q1):} $y=\tfrac12(\cos x-\sin x)+Ce^x$ \emph{is} a sol of $y'=y-\cos x$: $y'=\tfrac12(-\sin x-\cos x)+Ce^x=y-\cos x$.

\subsection{Finite-difference table drill (VII-1)}
Given $x:0.2,0.3,0.4$; $f:-0.3341,-0.1769,0.0138$ ($h=0.1$). Fwd: at $0.2$: $\tfrac{-0.1769+0.3341}{0.1}=1.572$; at $0.3$: $1.907$. Bwd = same shifted ($1.572$ at $0.3$, $1.907$ at $0.4$). Central at $0.3$: $\tfrac{0.0138+0.3341}{0.2}=1.7395=\tfrac12(1.907+1.572)$. 2nd: $\tfrac{0.0138-2(-0.1769)-0.3341}{0.01}=3.35$.

\subsection{Rectangle lattice nodes (Quiz 7)}
Given $x_k=\tfrac{3\ln(1+k/5)}{\ln2}-1$, $k=0..5$. Sample: $x_1=\tfrac{3\ln1.2}{\ln2}-1\approx-0.21090$ (since $3\ln1.2\approx0.54696$, $\ln2\approx0.69315$);\\
$x_0=\tfrac{3\ln1}{\ln2}-1=-1$; $x_5=\tfrac{3\ln2}{\ln2}-1=2$.\\
Full list: $-1.00000$, $-0.21090$, $0.45628$, $1.03422$, $1.54399$, $2.00000$.
$x_i^*=0.2x_{i-1}+0.8x_i$: $-0.36872$, $0.32285$, $0.91863$, $1.44204$, $1.90880$.\\
Widths $h_i=x_i-x_{i-1}$: $0.78910,0.66718,0.57794,0.50977,0.45601$; heights $f(x_i^*)$ ($f=x^3-1.25x^2$): $-0.22007,-0.09663,-0.27964,0.39934,2.40035$; e.g. $R_1=0.78910\cdot(-0.22007)\approx-0.17366$, total $R=\sum h_if(x_i^*)\approx0.89843$.\\
Rect $i$ has width $h_i=x_i-x_{i-1}$, height $f(x_i^*)$ ($f=x^3-1.25x^2$). \wk{How to draw (Pract.5):} (1) mark the $n{+}1$ nodes on the $x$-axis; (2) for each subinterval compute $x_i^*$ via the $\alpha$-formula (any $\alpha$); (3) go up to the curve at $f(x_i^*)$ — that height is the bar top; (4) draw the bar over $[x_{i-1},x_i]$ (works below the axis too: height negative $\Rightarrow$ signed area).

\subsection{Accuracy table $\int_0^1x^2dx$ (Lec17, $h=0.25$)}
Nodes $0,\tfrac14,\tfrac12,\tfrac34,1$; $f(x)=x^2$ gives $0,\tfrac1{16},\tfrac14,\tfrac9{16},1$.
Left $=\tfrac{17}{64}$? \hl{No:} $f_3=(\tfrac34)^2=\tfrac9{16}$, not $\tfrac34$. Correct left uses $f_0..f_3$: $=\tfrac14(0+\tfrac1{16}+\tfrac14+\tfrac9{16})=\tfrac14(\tfrac{14}{16})$\\
$=\tfrac7{32}\approx0.21875$, err $=|0.33333-0.21875|\approx0.11458$.\\
Mid: midpoints $\tfrac18,\tfrac38,\tfrac58,\tfrac78$, $f^*: \tfrac1{64},\tfrac9{64},\tfrac{25}{64},\tfrac{49}{64}$; $\tilde I=\tfrac14\cdot\tfrac{84}{64}\approx0.32813$ (err $0.00521$).\\
Right $=\tfrac14(\tfrac1{16}+\tfrac14+\tfrac9{16}+1)=\tfrac{15}{32}\approx0.46875$ (err $0.13542$).\\
Trap $=\tfrac{11}{32}\approx0.34375$ (err $0.01042$).
\wk{Derivation idea (Th.10):} Taylor + integrate: local $|E_i|\le M_1h_i^2/2$, sum $\Rightarrow|E|\le\tfrac{M_1(b-a)}2h=O(h)$; midpoint: odd term vanishes $\Rightarrow O(h^2)$.

\subsection{Simpson $\int_1^2e^{x^2}dx$ (Lec19, $n=4$)}
Given $h=(2-1)/4=0.25$, nodes $x_i=1+0.25i$: $1,\tfrac54,\tfrac32,\tfrac74,2$.
\wk{Step 1} — evaluate $f(x)=e^{x^2}$ (square $x$ first, then exponentiate): $f_0=e^{1}\approx2.71828$; $f_1=e^{(5/4)^2}=e^{1.5625}\approx4.77073$; $f_2=e^{(3/2)^2}=e^{2.25}\approx9.48774$; $f_3=e^{(7/4)^2}=e^{3.0625}\approx21.38094$; $f_4=e^{4}\approx54.59815$.\\
$I=\tfrac{0.25}3(2.71828+4(26.15167)+2(9.48774)+54.59815)\approx15.0749$.\\
\wk{Trap check:} two double-intervals $T_1=h(f_0{+}f_2)=0.25(2.71828+9.48774)=0.25(12.20602)\approx3.0515$, $T_2=h(f_2{+}f_4)=0.25(9.48774+54.59815)=0.25(64.08589)\approx16.0215$, total $\approx19.0730\to19.07$ — same ballpark $\Rightarrow$ likely right.\\
\wk{Simpson sketch (Quiz8 Q1):} $[0,240]$, $h=40$, 3 arcs (one per pair); mark nodes, each arc through its 3 points.

\subsection{Pseudoinverse $\to$ 1D formulas (HW09)}
$X^TX\bar a=X^T\bar y$ with $X=[1\ x_i]$ gives $\big(\begin{smallmatrix}n{+}1&\sum x_i\\\sum x_i&\sum x_i^2\end{smallmatrix}\big)\big(\begin{smallmatrix}a_0\\a_1\end{smallmatrix}\big)=\big(\begin{smallmatrix}\sum y_i\\\sum x_iy_i\end{smallmatrix}\big)$; dividing by $n{+}1$ and solving (Cramer) returns\\
$a_1=\tfrac{\overline{xy}-\bar x\bar y}{\overline{x^2}-\bar x^2}$, $a_0=\bar y-a_1\bar x$.
Cramer in full: $\Delta=(n{+}1)\sum x_i^2-(\sum x_i)^2$; $a_0=\tfrac{(\sum y_i)(\sum x_i^2)-(\sum x_i)(\sum x_iy_i)}{\Delta}$; $a_1=\tfrac{(n{+}1)(\sum x_iy_i)-(\sum x_i)(\sum y_i)}{\Delta}$; dividing top and bottom by $(n{+}1)^2$ gives the averaged formulas.
\ex{HW09 Q2 setup ($m{=}2,n{=}2$, 3 pts):} $X$ is $3\times3$ (first col of 1's, then $x_1^{(i)},x_2^{(i)}$ rows); $X^+=(X^TX)^{-1}X^T$ is $3\times3$. \ex{Q3:} $X^+X=I$ always; $XX^+=I$ only if $X$ square nonsingular.

\subsection{Euler $y'=y+x$ (Lec20, corrected)}
$(x_0,y_0)=(0,1)$, $h=0.25$: $y_1=1+0.25(1)=1.25$; $y_2=1.25+0.25(1.5)=1.625$;\\
$y_3=1.625+0.25(2.125)=2.15625$. (Compiled $1.5625/2.078125$ used $f_1=1.25$ instead of $0.25+1.25=1.5$ — knock-on slip.)\\
Exact $y^*=2e^x-x-1$: $y(0.5)\approx1.79744$, $y(0.75)\approx2.48400$ — Euler visibly below.
\ex{Midpoint on same ODE:} $(0.125,1.125)\to y_1=1+0.25(1.25)=1.3125$ vs exact $1.31805$ — much closer.\\
\hl{IVP slip (same ODE, HW pattern):} $y'=y+x$, $(1,4)$: $4=Ce-2\Rightarrow C=6/e\approx2.207$ (compiled $2/e$ forgot $+4$: check $y(1)=6-2=4$ $\checkmark$).

\subsection{Trapezoid, non-uniform lattice}
Given nodes $0,0.25,1$ $\Rightarrow$ widths $h_1=0.25-0=0.25$, $h_2=1-0.25=0.75$; $f=x^2$ gives $f_0=0$, $f_1=0.0625$, $f_2=1$.
$\tilde I=\tfrac12\big[(f_0{+}f_1)h_1+(f_1{+}f_2)h_2\big]$\\
$=\tfrac12\big[(0{+}0.0625)(0.25)+(0.0625{+}1)(0.75)\big]$\\
$=\tfrac12(0.015625+0.796875)=\tfrac12(0.8125)=0.40625$.\\
Exact $\tfrac13$; err $=|0.33333-0.40625|\approx0.07292$.

\subsection{Laplace expansion}
Given $A=\big(\begin{smallmatrix}1&2&3\\4&5&6\\7&8&10\end{smallmatrix}\big)$, expand along row 1.
\wk{Step 1} — cofactors: $\mathrm{cofac}(a_{11})=+(5\cdot10-6\cdot8)=50-48=2$; $\mathrm{cofac}(a_{12})=-(4\cdot10-7\cdot6)=-(40-42)=2$; $\mathrm{cofac}(a_{13})=+(4\cdot8-5\cdot7)=32-35=-3$.
\wk{Step 2} — assemble: $\det A=1(2)+2(2)+3(-3)=2+4-9=-3$ (matches LU $\det=1\cdot(-3)\cdot1$).\\
\ex{Row-swap + expand:} $D=\det\big(\begin{smallmatrix}5&6&7\\0&0&1\\8&9&10\end{smallmatrix}\big)$: swap rows 1,2 (flip sign), expand along the zero-heavy row 1:\\
$D=-(1\cdot(-1)^{1+3}(5\cdot9-6\cdot8))=-(-3)=3$.
\end{sloppypar}

\end{multicols}
\end{document}
